\documentclass[11pt]{article}
\usepackage[a4paper,margin=1in]{geometry}
\usepackage[T1]{fontenc}
\usepackage{lmodern}
\usepackage{booktabs}
\usepackage{longtable}
\usepackage{array}
\usepackage{graphicx}
\usepackage{hyperref}
\usepackage{enumitem}
\usepackage{xcolor}
\usepackage{setspace}

\hypersetup{
  colorlinks=true,
  linkcolor=blue,
  urlcolor=blue,
  citecolor=blue,
  pdftitle={Health Desert Scorer Nigeria Whitepaper},
  pdfauthor={Bashir},
  pdfsubject={LGA-level healthcare access decision support},
  pdfkeywords={Nigeria, health access, LGA, geospatial, machine learning}
}

\setlength{\parskip}{0.5em}
\setlength{\parindent}{0pt}
\onehalfspacing

\title{Health Desert Scorer (Nigeria)\\
\large A Decision-Support System for LGA-Level Healthcare Access Barriers}
\author{Bashir \\ \texttt{www.bashir.bio}}
\date{Draft v0.1 - \today}

\begin{document}
\maketitle

\begin{abstract}
The Health Desert Scorer is a Nigeria-focused decision-support application that identifies Local Government Areas (LGAs) with higher relative barriers to healthcare access. The system combines survey outcomes, facility distribution, distance proxies, and network connectivity signals to generate a comparative planning score. The tool is designed for prioritization and program planning by NGOs, clinics, policymakers, and research teams. It is explicitly not a diagnostic system and should not be used for individual-level prediction.
\end{abstract}

\tableofcontents
\newpage

\section{Executive Summary}
Healthcare access constraints are often driven by distance, facility scarcity, and weak connectivity, not only by clinical factors. This project provides an LGA-level prioritization layer for Nigeria so partners can:
\begin{itemize}[leftmargin=1.5em]
  \item identify highest-need LGAs with transparent criteria,
  \item compare risk dimensions (all-risk, child mortality, facility access, connectivity, 5km coverage),
  \item export planning-friendly artifacts for field operations and policy discussions.
\end{itemize}

The current model card reports:
\begin{itemize}[leftmargin=1.5em]
  \item Accuracy: 0.82
  \item Precision: 0.79
  \item Recall: 0.84
  \item F1: 0.81
  \item ROC-AUC: 0.88
\end{itemize}

\section{Problem Statement}
Many allocation decisions in public health are made with limited spatial synthesis across outcomes, services, and infrastructure. Existing dashboards can be difficult for non-technical users to operationalize into concrete outreach or deployment steps.

The Health Desert Scorer addresses this gap with a lightweight and shareable interface that ranks LGAs by relative access barriers and supports rapid planning conversations.

\section{System Overview}
\subsection{Product Scope}
The system is a map-first web application with:
\begin{itemize}[leftmargin=1.5em]
  \item state and year filtering,
  \item LGA search and comparison workflow,
  \item multiple focus modes for barrier-specific prioritization,
  \item downloadable outputs for meetings and offline review.
\end{itemize}

\subsection{Intended Users}
\begin{itemize}[leftmargin=1.5em]
  \item NGOs and implementation partners: intervention targeting and campaign planning.
  \item Primary care and clinic operators: outreach and catchment prioritization.
  \item Government stakeholders: transparent resource prioritization discussions.
  \item Researchers: hypothesis generation and sensitivity analysis support.
\end{itemize}

\subsection{Out of Scope}
\begin{itemize}[leftmargin=1.5em]
  \item Clinical diagnosis
  \item Individual-level risk prediction
  \item Use as sole basis for high-stakes funding decisions
\end{itemize}

\section{Data Inputs and Governance}
\subsection{Primary Data Inputs}
\begin{longtable}{p{3.0cm}p{2.8cm}p{2.0cm}p{6.0cm}}
\toprule
\textbf{Source} & \textbf{Dataset} & \textbf{Year(s)} & \textbf{Role in System} \\
\midrule
DHS & Survey outcomes & 2013, 2018 & Outcome and risk proxy signals at aggregated level \\
NHFR & Facility registry & 2020 & Facility density and access proxy features \\
WorldPop & Population estimates & 2020 & Population-normalized metrics (e.g., facilities per 10k) \\
OpenCellID & Connectivity proxy & 2019 & Tower density and mobile connectivity context \\
\bottomrule
\end{longtable}

\subsection{Data Protection and Ethics}
The repository stores aggregated outputs only. Restricted microdata and sensitive coordinates are not committed. This design supports safer sharing and reproducibility for collaborative planning.

\section{Methodology}
\subsection{Feature Engineering}
The current model family relies on core LGA-level features, including:
\begin{itemize}[leftmargin=1.5em]
  \item under-5 mortality proxy,
  \item facilities per 10k population,
  \item average distance to facilities,
  \item towers per 10k population,
  \item 5km coverage proxy,
  \item additional normalized contextual terms.
\end{itemize}

\subsection{Modeling}
The production model card describes a gradient boosted tree model (versioned artifact \texttt{risk\_model\_v1.2}) with 5-fold cross-validation stratified by state.

\subsection{Evaluation Metrics}
\begin{table}[h]
\centering
\begin{tabular}{lc}
\toprule
\textbf{Metric} & \textbf{Value} \\
\midrule
Accuracy & 0.82 \\
Precision & 0.79 \\
Recall & 0.84 \\
F1 Score & 0.81 \\
ROC-AUC & 0.88 \\
\bottomrule
\end{tabular}
\caption{Current performance snapshot from \texttt{models/risk\_model\_v1.2/metrics.json}.}
\end{table}

\subsection{Interpretability}
The application supports feature-attribution views for selected LGAs using SHAP-like local explanation payloads where available. Interpretability is included to improve trust in planning conversations, not to imply causal inference.

\section{User Workflows}
\subsection{NGO Workflow}
\begin{enumerate}[leftmargin=1.5em]
  \item Filter to state and year.
  \item Rank LGAs by selected focus (e.g., facility access).
  \item Export top-priority list and map snapshot.
  \item Validate with field and partner constraints.
\end{enumerate}

\subsection{Clinic and Frontline Workflow}
\begin{enumerate}[leftmargin=1.5em]
  \item Select service area state.
  \item Compare 2--4 LGAs for outreach sequencing.
  \item Use barrier-specific recommendations as planning prompts.
\end{enumerate}

\subsection{Research Workflow}
\begin{enumerate}[leftmargin=1.5em]
  \item Export model and feature outputs.
  \item Reproduce scores and inspect feature importances.
  \item Run sensitivity checks by year, source coverage, and thresholds.
\end{enumerate}

\section{Limitations}
Known limitations include:
\begin{itemize}[leftmargin=1.5em]
  \item no direct encoding of conflict and insecurity constraints,
  \item no seasonal road accessibility model in current release,
  \item care quality and staffing adequacy are not directly modeled,
  \item source completeness can vary across states and years.
\end{itemize}

\section{Deployment and Operations}
\subsection{Current Deployment}
The live system is delivered as a Streamlit-hosted app with an embedded custom HTML/JavaScript frontend.

\subsection{Operational Readiness Checklist}
\begin{itemize}[leftmargin=1.5em]
  \item versioned model artifact and model card published,
  \item data freshness metadata visible in UI/export artifacts,
  \item per-LGA quality flags surfaced to end-users,
  \item pilot governance process for partner sign-off and review.
\end{itemize}

\section{Pilot Design for External Partners}
\subsection{Proposed Pilot Structure (90 Days)}
\begin{itemize}[leftmargin=1.5em]
  \item Week 1--2: stakeholder alignment and success metric definition.
  \item Week 3--6: deployment, onboarding, and scenario testing.
  \item Week 7--10: live planning use during at least one campaign cycle.
  \item Week 11--12: evaluation, refinement, and scale decision memo.
\end{itemize}

\subsection{Pilot Success Metrics}
\begin{itemize}[leftmargin=1.5em]
  \item reduction in planning cycle time,
  \item percentage of outreach plans supported by explicit spatial evidence,
  \item adoption rate among intended user roles,
  \item qualitative trust and usability scores.
\end{itemize}

\section{Commercialization Snapshot}
The recommended commercialization pathway is service-led software:
\begin{itemize}[leftmargin=1.5em]
  \item paid pilot implementations,
  \item annual subscriptions for hosted access and updates,
  \item paid custom analytics packages for donor and research workflows.
\end{itemize}

\section{Conclusion}
The Health Desert Scorer is best positioned as a practical planning layer between raw health data systems and field-level decision-making. Its value is in speed, usability, and transparent prioritization. With structured pilots, role-based onboarding, and documented governance, the project can evolve into a durable product for NGOs, clinic networks, and research partners.

\appendix
\section{Artifacts to Include Before Public Submission}
\begin{itemize}[leftmargin=1.5em]
  \item Full model card and changelog
  \item Data dictionary and source refresh dates
  \item Validation protocol and pilot report template
  \item Responsible use policy and escalation flow
\end{itemize}

\section{Build Instructions}
From repository root:
\begin{verbatim}
latexmk -pdf -output-directory=docs/whitepaper docs/whitepaper/health_desert_whitepaper.tex
\end{verbatim}

\end{document}
